%!TEX root = ../../main.tex
\documentclass[tikz]{standalone}

\usepackage{amsmath, amssymb}
\usepackage{tikz}
\usetikzlibrary{arrows.meta, calc, positioning, shapes.geometric, angles, quotes, backgrounds}
\usepackage{subcaption}

\definecolor{academicBlue}{RGB}{0, 85, 164}
\definecolor{academicRed}{RGB}{200, 30, 40}
\definecolor{marginGray}{RGB}{220, 220, 220}

\begin{document}
	\begin{figure}[H]
	    \centering
	    \tikzset{
	        >=Stealth,
	        feature_point/.style={circle, inner sep=1.5pt, fill=black!70},
	        class1_point/.style={feature_point, fill=academicBlue},
	        class2_point/.style={feature_point, fill=academicRed},
	        weight_vec/.style={->, thick, shorten >=1pt},
	        boundary_line/.style={dashed, thick, black!60},
	        margin_area/.style={fill=marginGray, opacity=0.5},
	        axis_style/.style={->, thin, black!40}
	    }

	    % --- Softmax ---
	    \subfloat[Standard Softmax with Cross-Entropy loss]{%
	        \label{fig:softmax_geom}%
	        \begin{tikzpicture}[scale=2.5]
	            \coordinate (O) at (0,0);
	            \draw[axis_style] (-0.2,0) -- (1.8,0);
	            \draw[axis_style] (0,-0.2) -- (0,1.8);
            
	            \def\angWone{30}
	            \def\angWtwo{130}
	            \pgfmathsetmacro{\angBoundary}{(\angWone + \angWtwo)/2}

	            \coordinate (W1_tip) at (\angWone:1.4);
	            \coordinate (W2_tip) at (\angWtwo:1.2);

	            \draw[weight_vec, academicBlue] (O) -- (W1_tip) node[right] {$\mathbf{W}_1$};
	            \draw[weight_vec, academicRed] (O) -- (W2_tip) node[left] {$\mathbf{W}_2$};
	            \draw[boundary_line] (O) -- (\angBoundary:1.8) node[above right, font=\footnotesize] {Decision Boundary};

	            \coordinate (X_sample) at (\angWone+15:1.1);
	            \draw[->, thin, black!80] (O) -- (X_sample) node[midway, below, font=\footnotesize] {$\mathbf{x}_i$};
            
	            \pic [draw, angle radius=0.5cm, "$\theta_{1,i}$", font=\footnotesize, angle eccentricity=1.3] {angle = X_sample--O--W1_tip};

	            \foreach \r/\a in {1.0/20, 1.3/35, 0.8/25, 1.5/40, 1.1/15, 0.9/30, 1.2/10}
	                \node[class1_point] at (\a:\r) {};
	            \foreach \r/\a in {1.1/120, 0.9/135, 1.4/125, 1.0/140, 1.2/115, 0.8/130, 1.3/145}
	                \node[class2_point] at (\a:\r) {};
	        \end{tikzpicture}%
	    }
	    \hfil
	    % --- ArcFace ---
	    \subfloat[ArcFace Loss]{%
	        \label{fig:arcface_geom}%
	        \begin{tikzpicture}[scale=1.7]
	            \coordinate (O) at (0,0);
	            \def\radius{1.5}
	            \draw[thick, black!80] (O) circle (\radius);
	            \node at (115:\radius+0.3) [font=\footnotesize] {Hypersphere};

	            \def\angWone{30}
	            \def\angWtwo{130}
	            \def\marginDeg{25}
	            \pgfmathsetmacro{\angBoundary}{(\angWone + \angWtwo)/2}
	            \pgfmathsetmacro{\angMarginOne}{\angBoundary - \marginDeg}
	            \pgfmathsetmacro{\angMarginTwo}{\angBoundary + \marginDeg}

	            \coordinate (W1_pt) at (\angWone:\radius);
	            \coordinate (W2_pt) at (\angWtwo:\radius+0.05);
	            \coordinate (M1_pt) at (\angMarginOne:\radius);
	            \coordinate (M2_pt) at (\angMarginTwo:\radius);

	            \begin{scope}[on background layer]
	                \fill[marginGray] (O) -- (M1_pt) arc (\angMarginOne:\angMarginTwo:\radius) -- cycle;
	            \end{scope}
            
	            \draw[boundary_line, thin] (O) -- (M1_pt);
	            \draw[boundary_line, thin] (O) -- (M2_pt);
	            \draw[dashed, thin, black!30] (O) -- (\angBoundary:\radius);

	            \draw[weight_vec, academicBlue] (O) -- (W1_pt) node[right] {$\mathbf{W}_1$};
	            \draw[weight_vec, academicRed] (O) -- (W2_pt) node[left] {$\mathbf{W}_2$};

	            \pic [draw, <->, angle radius=0.8*\radius cm, "$m$", font=\small, angle eccentricity=1.15, academicBlue] {angle = W1_pt--O--M1_pt};
	            \pic [draw, <->, angle radius=0.8*\radius cm, "$m$", font=\small, angle eccentricity=1.15, academicRed] {angle = M2_pt--O--W2_pt};

	            \foreach \a in {28, 32, 25, 35, 30, 22, 38}
	                \node[class1_point] at (\a:\radius) {};
	            \foreach \a in {128, 132, 125, 135, 130, 122, 138}
	                \node[class2_point] at (\a:\radius) {};

	            \node at (\angBoundary:0.6*\radius) [font=\footnotesize, align=center, rotate=\angBoundary-90] {Margin\\Gap};
	        \end{tikzpicture}%
	    }

	    \caption{Geometric interpretation of Softmax vs. ArcFace. (a) Softmax creates linear boundaries without intra-class compactness. (b) ArcFace projects features onto a hypersphere and enforces an angular margin $m$, creating a clear gap between classes.}
	    \label{fig:arcface_geometric_intuition}
	\end{figure}
\end{document}
